\chapter{Дополнительные формулы и таблицы для главы 3 и части главы 4}

% =========================================================
% Основа приложения:
% - Приложение B/C из предоставленного пользователем файла
% - рабочий план диссертации
%
% По рабочему плану сюда при необходимости можно дополнительно включить:
% - громоздкие формулы для линейных преобразований;
% - дополнительные таблицы коэффициентов для главы 3;
% - вспомогательные вычислительные материалы по прикладному алгоритму
%   из главы 4, если их включение в основной текст ухудшает связность.
% =========================================================

\textit{Примечание.}
В текущей версии приложения собраны табличные коэффициенты
$R_{i\Omega}$-преобразования и формулы для восстановления коэффициентов
аппроксимации по эталонным значениям сопряжённых переменных при двух
соотношениях радиусов. При дальнейшей доработке сюда могут быть
добавлены другие громоздкие формулы, таблицы коэффициентов и
вспомогательные вычислительные материалы для линейных преобразований и
прикладного алгоритма главы 4.

\begin{table}[htbp]
\centering
\caption{Коэффициенты аппроксимации $R_{i\Omega}$-преобразования}
\label{tab:appendixC_RiOmega}
\begin{tabular}{|c|c|c|}
\hline
$N$ & $C^{i}_{0}$ & $C^{i}_{1}$ \\
\hline
0 & $-4.459247\mathrm{E}{-2}$ & $1.208435\mathrm{E}{-3}$ \\
1 & $-2.516674\mathrm{E}{-1}$ & $-1.425563$ \\
2 & $3.938937\mathrm{E}{-2}$ & --- \\
\hline
\end{tabular}
\end{table}

\noindent
Коэффициенты $C_i^{\kappa}(\lambda,\kappa_1,\kappa_2)$ рассчитываются на
основе эталонных значений сопряжённых переменных для двух соотношений
радиусов.

\begin{equation}
C_0^{\kappa}(\lambda,\kappa_1,\kappa_2)=
\frac{\kappa_1^{3/4}\kappa_2^{3/4}}{\kappa_2^{3/4}-\kappa_1^{3/4}}
\left(
\kappa_1^{1/4}\frac{p_{rx}(\lambda,\kappa_1)}{\log\kappa_1}-
\kappa_2^{1/4}\frac{p_{rx}(\lambda,\kappa_2)}{\log\kappa_2}
\right)
\tag{123}
\end{equation}

\begin{equation}
C_1^{\kappa}(\lambda,\kappa_1,\kappa_2)=
\frac{1}{\kappa_1^{3/4}-\kappa_2^{3/4}}
\left(
\kappa_1\frac{p_{rx}(\lambda,\kappa_1)}{\log\kappa_1}-
\kappa_2\frac{p_{rx}(\lambda,\kappa_2)}{\log\kappa_2}
\right)
\tag{124}
\end{equation}

\begin{equation}
C_2^{\kappa}(\lambda,\kappa_1,\kappa_2)=
\frac{\kappa_1^{1/4}\kappa_2^{1/4}}{\kappa_2^{1/4}-\kappa_1^{1/4}}
\left(
\kappa_1^{3/4}\frac{p_{ry}(\lambda,\kappa_1)}{\log\kappa_1}-
\kappa_2^{3/4}\frac{p_{ry}(\lambda,\kappa_2)}{\log\kappa_2}
\right)
\tag{125}
\end{equation}

\begin{equation}
C_3^{\kappa}(\lambda,\kappa_1,\kappa_2)=
\frac{1}{\kappa_1^{1/4}-\kappa_2^{1/4}}
\left(
\kappa_1\frac{p_{ry}(\lambda,\kappa_1)}{\log\kappa_1}-
\kappa_2\frac{p_{ry}(\lambda,\kappa_2)}{\log\kappa_2}
\right)
\tag{126}
\end{equation}

\begin{equation}
C_4^{\kappa}(\lambda,\kappa_1,\kappa_2)=
\frac{\kappa_1^{1/4}\kappa_2^{1/4}}{\kappa_2^{1/4}-\kappa_1^{1/4}}
\left(
\kappa_1^{3/4}\frac{p_{vx}(\lambda,\kappa_1)}{\log\kappa_1}-
\kappa_2^{3/4}\frac{p_{vx}(\lambda,\kappa_2)}{\log\kappa_2}
\right)
\tag{127}
\end{equation}

\begin{equation}
C_5^{\kappa}(\lambda,\kappa_1,\kappa_2)=
\frac{1}{\kappa_1^{1/4}-\kappa_2^{1/4}}
\left(
\kappa_1\frac{p_{vx}(\lambda,\kappa_1)}{\log\kappa_1}-
\kappa_2\frac{p_{vx}(\lambda,\kappa_2)}{\log\kappa_2}
\right)
\tag{128}
\end{equation}

\begin{equation}
C_6^{\kappa}(\lambda,\kappa_1,\kappa_2)=
\frac{\kappa_1^{3/4}\kappa_2^{3/4}}{\kappa_2^{3/4}-\kappa_1^{3/4}}
\left(
\kappa_1^{1/4}\frac{p_{vy}(\lambda,\kappa_1)}{\log\kappa_1}-
\kappa_2^{1/4}\frac{p_{vy}(\lambda,\kappa_2)}{\log\kappa_2}
\right)
\tag{129}
\end{equation}

\begin{equation}
C_7^{\kappa}(\lambda,\kappa_1,\kappa_2)=
\frac{1}{\kappa_1^{3/4}-\kappa_2^{3/4}}
\left(
\kappa_1\frac{p_{vy}(\lambda,\kappa_1)}{\log\kappa_1}-
\kappa_2\frac{p_{vy}(\lambda,\kappa_2)}{\log\kappa_2}
\right)
\tag{130}
\end{equation}
